\documentclass{beamer}

\usepackage{amsmath, amssymb}
\usepackage{xcolor}


\title{Exercise 2.2: Helmholtz Problem \\ Pollution Effect}
\author{}
\date{}

\begin{document}

\frame{\titlepage}

%-------------------------------------------------
\begin{frame}{Problem Statement}
We consider the Helmholtz problem on the unit ball $\Omega \subset \mathbb{R}^2$:
\[
(-k^{-2}\Delta - 1)u = 0 \quad \text{in } \Omega
\]
\[
\frac{\partial u}{\partial n} - iku = g \quad \text{on } \partial \Omega
\]

with
\[
g(x) = ik(x_1 - 1)e^{ikx_1}.
\]

\bigskip

\textbf{Goals:}
\begin{itemize}
    \item Show that $u = e^{ikx_1}$ is a solution
    \item Prove that $u$ satisfies the variational formulation
\end{itemize}
\end{frame}

%-------------------------------------------------
\begin{frame}{Task 1: Verifying the PDE}
Ansatz:
\[
u(x) = e^{ikx_1}
\]

\textbf{Compute derivatives:}
\[
\nabla u = (ik e^{ikx_1}, 0)^T, \quad
\Delta u = -k^2 e^{ikx_1}
\]

\bigskip

Insert into PDE:
\begin{align*}
    (-k^{-2}\Delta -1)u
&=-k^{-2}\Delta u-u \\
&= -k^{-2}(-k^2e^{ikx_1}) -e^{ikx_1}\\
&= (1-1)e^{ikx_1} = 0
\end{align*}

\bigskip

\textbf{$\Rightarrow$ PDE is satisfied}
\end{frame}

%-------------------------------------------------
\begin{frame}{Task 1: Boundary Condition}
Normal derivative:
\[
\frac{\partial u}{\partial n} = \nabla u \cdot n
= ik e^{ikx_1} n_1
\]

Boundary expression:
\[
\frac{\partial u}{\partial n} - iku
= ik e^{ikx_1}(n_1 - 1)
\]

On $\partial \Omega$:
\[
n = x \Rightarrow n_1 = x_1
\]

\bigskip

\[
\Rightarrow g(x) = ik(x_1 -1)e^{ikx_1}
\]

\bigskip

\textbf{$\Rightarrow$ Boundary condition is satisfied}
\end{frame}

%-------------------------------------------------
\begin{frame}{Task 2: Variational Formulation}
We want to show:
\[
a_k(u,v) = L(v) \quad \forall v \in H^1(\Omega)
\]

with
\[
a_k(u,v) = \int_\Omega \left(k^{-2} \nabla u \cdot \nabla \overline{v} - u\overline{v}\right) dx
- ik^{-1} \int_{\partial \Omega} u\overline{v} \, ds
\]

\[
L(v) = k^{-2} \int_{\partial \Omega} g \overline{v} \, ds
\]

\end{frame}

%-------------------------------------------------
\begin{frame}{Green's Identity}
We use:
\[
\int_\Omega \nabla u \cdot \nabla \overline{v} dx
= -\int_\Omega (\Delta u)\overline{v} dx
+ \int_{\partial \Omega} \frac{\partial u}{\partial n} \overline{v} ds
\]

Insert into $a_k(u,v)$:

\begin{align*}
a_k(u,v) = &\int_\Omega \left(k^{-2} \nabla u \cdot \nabla \overline{v} - u\overline{v}\right) dx
- ik^{-1} \int_{\partial \Omega} u\overline{v} \, ds\\
= &-k^{-2} \int_\Omega (\Delta u)\overline{v} dx
+ k^{-2} \int_{\partial \Omega} \frac{\partial u}{\partial n} \overline{v} ds \\
&- \int_\Omega u\overline{v} dx
- ik^{-1} \int_{\partial \Omega} u\overline{v} ds  
\end{align*}
\end{frame}


\begin{frame}{}
Insert \[u(x) = e^{ikx_1}\]
and using the normal derivative on the boundary:
\[
\frac{\partial u}{\partial n}(x) = ik e^{ikx_1} x_1
\]  and\[\Delta u = -k^2e^{ikx_1}\]

\begin{align*}
a_k(u,v)
=& \textcolor{blue}{-k^{-2} \int_\Omega (-k^2e^{ikx_1})\overline{v} dx}
+ k^{-2} \int_{\partial \Omega} ik e^{ikx_1}\textcolor{red}{\mathbf{ x_1}} \overline{v} ds
\\& \textcolor{blue}{- \int_\Omega e^{ikx_1}\overline{v} dx}
- ik^{-1} \int_{\partial \Omega} e^{ikx_1}\overline{v} ds\\
&= ik^{-1} \int_{\partial \Omega} (x_1 -1) e^{ikx_1}\overline{v} ds\\
&= ik^{-2} \int_{\partial \Omega} g\overline{v} ds\\
&= L(v)
\end{align*}
\end{frame}


\begin{frame}{Exercise 6}
    \includegraphics[width=1.05\textwidth]{Exercise_6.PNG}
    \includegraphics[width=1.05\textwidth]{Code_Exercise_6.PNG}

\end{frame}

\begin{frame}{Exercise 6}
    \begin{columns}
    \begin{column}[c]{0.48\textwidth}
        \hbox{}
        
        \vspace{1.em}
        \begin{minipage}{\textwidth}
            \includegraphics[width=.99\textwidth]{Grid_Exercise_6.PNG}
        \end{minipage}
        \end{column}
        \begin{column}[c]{0.48\textwidth}
        \begin{minipage}{\textwidth}        
            \includegraphics[width=.99\textwidth]{Solution_Exercise_6.PNG}
        \end{minipage}
    \end{column}
    \end{columns}
\end{frame}

\end{document}